\input{../tex-inputs/latex-leseansicht-vorspann}

\section[Arthur Schnitzler: Widmungsexemplar Frau Bertha Garlan für Theodor von Sosnosky, {[}24.?{]} 4. 1901]{L04238 Arthur Schnitzler: Widmungsexemplar Frau Bertha Garlan für Theodor von
               Sosnosky, {[}24.?{]} 4. 1901}
\nopagebreak\mylabel{L04238v}
\rehead{ }\normalsize\beginnumbering\briefempfaengerindex{Sosnosky, Theodor von@\textsc{Sosnosky, Theodor von}!zzzSchnitzler, Arthur@\emph{von Arthur Schnitzler}!1901-04-243@{{[}24.?{]} 4. 1901}|(be}
\toendnotes[C]{\smallbreak\pagebreak[2]}
\correspDesc{Versand  durch Arthur Schnitzler am [24.?] 4. 1901 in Wien
\newline{}Erhalt  durch Theodor von Sosnosky im Zeitraum [25. 4. 1901 – 29. 4. 1901?] \textbf{Ort fehlend} }\toendnotes[C]{\smallbreak}
\Standort{Privatbesitz, Nico Dostal, \emph{ohne Signatur}.}
\physDesc{Widmung am Vorsatzblatt, 88 Zeichen
\newline{}Handschrift: schwarze Tinte, deutsche Kurrent
\newline{}Ordnung: Vermerk »Ger« }\toendnotes[C]{\smallbreak}
\pstart
           \noindent{}\centering{}{\pb}Herrn Theodor von Sosnosky in aufrichtiger Hochſchätzung\pend
           \pstart \spacefill\mbox{ArthSchnitzler}\pend{}
\pstart
           \noindent{}Wien\oindex{Wien@\textbf{Wien}, \emph{Verwaltungsgebiet}|pw}, \label{K_L04238-1v}\edtext{April 901}{\lemma{\textnormal{\emph{April 901}}}\Cendnote{\textnormal{Die Datierung erfolgt in Abgleich zum
                        Widmungsexemplar für Bahr\pwindex{Bahr, Hermann 19.\,7.\,1863 Linz – 15.\,1.\,1934 München@\textsc{Bahr, Hermann} (19.\,7.\,1863 Linz – 15.\,1.\,1934 München), \emph{Schriftsteller, Kritiker}|pwk}, für das
                        sich dieser am XXXX Auszeichnungsfehler: Dokument L01115 nicht gefunden bedankt hat.}}}\label{K_L04238-1}.\pend
           {\vspace{1\baselineskip}}\selectlanguage{ngerman}\vspace{1em}{\vspace{1\baselineskip}}
\pstart
           \centering{}{\pb}\textcolor{gray}{\textbf{Frau Bertha Garlan\pwindex{Schnitzler, Arthur 15.\,5.\,1862 Wien – 21.\,10.\,1931 ebd.@\textsc{Schnitzler, Arthur} (15.\,5.\,1862 Wien – 21.\,10.\,1931 ebd.), \emph{Schriftsteller, Mediziner}!Frau Bertha Garlan. Roman@\strich\emph{Frau Bertha Garlan. Roman}|pw}}}\pend
           
\pstart
           \centering{}\textcolor{gray}{\textbf{\so{Roman}}}\pend
           
\pstart
           \centering{}\textcolor{gray}{\textbf{von}}\pend
           
\pstart
           \centering{}\textcolor{gray}{\textbf{Arthur Schnitzler}}\pend
           {\vspace{1\baselineskip}}
\pstart
           \centering{}\textcolor{gray}{\textbf{\textbf{Berlin}\oindex{Berlin@\textbf{Berlin}, \emph{Hauptstadt}|pw}}}\pend
           
\pstart
           \centering{}\textcolor{gray}{\textbf{\so{S. Fiſcher, Verlag}\orgindex{S. Fischer Verlag@S. Fischer Verlag|pw}}}\pend
           
\pstart
           \centering{}\textcolor{gray}{\textbf{1901}}\pend
           \selectlanguage{ngerman}\endnumbering\briefempfaengerindex{Sosnosky, Theodor von@\textsc{Sosnosky, Theodor von}!zzzSchnitzler, Arthur@\emph{von Arthur Schnitzler}!1901-04-243@{{[}24.?{]} 4. 1901}|)be}\mylabel{L04238h}
\begin{anhang}
\end{anhang}\newcommand{\dateiname}{L04238}\newcommand{\titel}{Arthur Schnitzler: Widmungsexemplar Frau Bertha Garlan für Theodor von Sosnosky, [24.?] 4. 1901}\newcommand{\editorInnen}{Herausgegeben von Jahnke, SelmaMüller, Martin Anton}\input{../tex-inputs/latex-leseansicht-abspann}
